\section{Related Work}
\label{sec:related}

This section positions Cstamp-PWAL within the landscape of crash recovery and
durability optimization techniques.  We discuss parallel WAL designs, dependency
tracking, timestamp-based recovery, and latency-aware durability protocols.

\subsection{Parallel Write-Ahead Logging}

Conventional centralized WAL (e.g., ARIES) serializes all log insertion and
durable acknowledgment through a single log stream.  Parallel WAL (P-WAL) designs
address this bottleneck by partitioning log records across multiple shards,
allowing concurrent appends from different workers.

However, P-WAL designs still differ significantly in how they handle durable
acknowledgment.  Worker-wait designs, where application threads block until their
log record is persisted, ensure tight coupling between commit and durability but
eliminate pipeline parallelism.  Asynchronous designs decouple worker threads
from flushing, but still face the question: when can a transaction be
acknowledged as durable?

\subsection{Global-Prefix Durability}

The simplest asynchronous durability protocol allocates a monotonically
increasing log sequence number (LSN) to each transaction and acknowledges a
transaction only when a global durable prefix reaches its LSN.  This rule is
safe but conservative: it forces transactions on fast shards to wait for slow
shards, even when there are no actual dependencies.  Such conservative waiting
accumulates durable-acknowledgment backlog and increases tail latency.

\subsection{Dependency-Based Logging and Recovery}

Several systems track dependencies to avoid conservative durable-prefix waiting.
Taurus~\cite{Taurus} proposes dependency-based logging, where recovery order is
determined by actual read-write dependencies rather than global LSN order.  This
approach reduces unnecessary ordering constraints.

However, Taurus targets distributed systems where log records may be on different
machines, and its design involves explicit dependency tracking at the logging
layer.  RFA (Remote Flush Avoidance)~\cite{RFA} applies similar ideas in
multi-node settings.  CPR (Checkpoint-based Parallel Recovery)~\cite{CPR}
optimizes recovery by grouping transactions into dependency-closed checkpoints,
reducing the state that must be replayed.

\subsection{Cstamp-PWAL: Integration with Timestamp-Based Engines}

Cstamp-PWAL builds on dependency-based recovery ideas, but targets single-node
main-memory transaction processing engines that already assign commit timestamps
(e.g., ERMIA).  Our key contributions beyond existing dependency-based logging
are:

\begin{enumerate}

\item \emph{Reuse of commit timestamp as logical LSN}: Existing WAL systems
separate commit timestamps (used for concurrency control) from recovery order
(defined by WAL LSN).  Cstamp-PWAL eliminates this duplication by using the
commit timestamp itself as the logical recovery order.  This simplification
avoids allocating a separate global WAL LSN on the hot path.

\item \emph{Dependency frontier derived from engine-level dependencies}:
Cstamp-PWAL computes the dependency frontier directly from read-write
dependencies already tracked by the transaction engine (e.g., ERMIA's own
validation mechanism).  This allows tight integration with the concurrency
control layer without adding a separate dependency-tracking abstraction.

\item \emph{Evaluation of durable-ack backlog and tail latency}: Most prior work
on dependency-based recovery focuses on correctness or throughput.  Cstamp-PWAL
explicitly measures and optimizes for durable-acknowledgment backlog and tail
latency, which are critical for latency-sensitive workloads.  The key insight
is that on sparse-dependency workloads, reducing backlog and tail latency is as
important as improving peak throughput.

\item \emph{Worker/flusher/committer pipeline}: Rather than having workers block
on either flush or dependency waits, Cstamp-PWAL separates these concerns into
three pipeline stages.  This design allows workers to continue executing while
independent transactions' log records are being flushed, without blocking on
conservative global-prefix waits.

\end{enumerate}

In summary, while dependency-based logging is not a new idea, Cstamp-PWAL
combines dependency tracking with timestamp-based engines and evaluates the
practical impact on durable-acknowledgment latency---a gap that existing
dependency-based logging papers do not thoroughly address.

\subsection{Epoch-Based and Timestamp-Based Durability}

SiloR~\cite{SiloR} extends the Silo transaction engine with crash recovery by
using epoch-based checkpoints and log truncation.  Silo's fixed-epoch design
simplifies recovery but limits flexibility in batching.

Cstamp-PWAL, by contrast, uses fine-grained commit timestamps (often assigned at
transaction commit time) rather than coarse-grained epochs, allowing finer-
grained control over durable-acknowledgment boundaries.

\subsection{Latency-Aware Durability}

Recent systems have begun to address the tail-latency impact of durability
protocols.  Cstamp-PWAL is motivated by similar concerns: conservative global-
prefix waiting not only hurts peak throughput but also causes committed
transactions to accumulate in memory, waiting for acknowledgment.  By replacing
global-prefix coupling with dependency-closed acknowledgment, Cstamp-PWAL
reduces both backlog and tail latency on sparse-dependency workloads.

